% Pursuing Stacks (cotes 134-2 to 134-8): archive cote 134-3, site page 48 (PDF page 49), typescript page 48.
% typescript. Draft compared: ps2.tex:2385-2457.
% AI-transcribed by Claude Opus 5.5 (claude-opus-5-5, adaptive thinking, effort high), Message Batches API, batch ps1 (msgbatch_01Rv5hCaNTtUt4JHx1yqoHPL), 2026-11-12; not yet checked by eye.
\scanpage{134-3/48}
\mspage{48}

Of course,  the element I  (playing the role of unit interval with
endpoints $e_o$,$e_1$) is by no means unique,  for instance it can be replaced
by any cartesian power $I^n$ $(n\geqslant 1)$,  provided it is in \over{A}{$A$}  -  or in the case
of A = \hand{$\Delta$}  we can take for I any \hand{$\Delta$}$_n$ $(n\geqslant 1)$ and for \hand{$\delta$}$_o$ ,\hand{$\delta$}$_1$ any two \uline{distinct}
maps from e=\hand{$\Delta$}$_o$ to \hand{$\Delta$}$_n$ , instead of the usual choice \xout{for} \hand{$\Delta$}$_1$ for I, and
corresponding \hand{$\delta$}$_o$ and \hand{$\delta$}$_1$ .  This high degree of arbitrar\over{y}{i}ness in the choice
of I should be no surprise,  this was already one striking feature in Quil-%
len's theory of model categories.

There remains the asphericity condition (T 2)  for the categories
$A_{/axb}$ ,  with  a,b \hand{$\in$} Ob A ,  which is somewhat subtler.  There is one very
evident way though to insure this, namely assuming

T 2')  If a,b are in A,  so is axb  i.e. A is stable under binary products.

In other words,  putting together  T 1' and T 2',  \xout{taking any categ} we
may \xout{restrict} take categories (essential small) \uline{stable under finite products}
When such a category satisfies the mild extra condition T 3') above,  it
is a test category!  This is already an impressive bunch of test categories
For instance take any category \inserted{C}\inserted[typed]{with} \xout{stable under} finite products and \xout{for which}
for which there exist I,$e_o$,$e_1$ as in T 3'  -  never mind whether \hand{C} is essen-%
tially small,  for instance any "non empty" topos will do (taking for I
the Lawvere element for instance,  if no simpler choice comes to \hand{m}ind).
Ta\hand{k}e any subcategory \inserted{A} (full or not) stable \xout{by}\inserted[typed]{under} finite products and containing
I and \hand{$\Delta$}$_o$ ,\hand{$\Delta$}$_1$ and where the  C-products are also A-products  -  for instan\xout{c}
ce this is OK if A is full.  Then A is a test category. \margin{$\lceil$} The simplest choi-%
ce her\over{e}{e},  the smallest in any case,  is to take the subcategory made up
with all cartesian powers $I^n$ $(n\geqslant 0)$.  We may take the full subcategory
made up with these elements,  but instead we may take,\xout{as only arrows} still
more economically,  only the arrows
\[ I^n \rightarrow I^m \]
whose \inserted{m} components  $I^n \rightarrow I$ are each,  either a projection $pr_i$ $(1\leqslant i\leqslant n)$ or
of the form \hand{$\delta$}$_i$\hand{$($}$p_{I^n}$\hand{$)$} $(i$\hand{$\in\{$}$0,1$\hand{$\}$}$)$.  The\xout{\unsure{s}} category thus obtained,  up to unique
isomorphism,  does not depend on the choice of (C,I,$e_o$,$e_1$),  visibly -
it is,  in a sense,  \uline{the} smallest test category satisfying the stronger
conditions T 1' to T 3'.  \xout{One} It's elements can be visualized as "standar\hand{d}
cubes"\over{\unsure{.}}{\unsure{:}} One convenient way to do so is to take C = (Ord) (category of
all ordered sets), I = \hand{$\Delta$}$_1$ = ( \hand{$\{$}0\hand{$\}$} \hand{$\rightarrow$} \hand{$\{$}1\hand{$\}$} ) ,  $e_o$ and $e_1$ as usual,  thus
A can be interpreted as a subcategory of (Ord),  but this embedding is not
full (there are maps of $I^2$ to I in (Ord)  which are not in A,  \unsure{i.e,} do not
"respect the cubical structure").  We may also take  C = (Spaces),  I =
unit interval \hand{$\subset$} $\underline{R}$ ,  $e_o$ and $e_1$ defined by the endpoints  O and  1 ,
